\protect \centerline {\protect \bf  LỜI GIẢI: ĐỀ THI HỌC KÌ 2017-2018-SỐ 1}
\begin{Soln}{1}
Bài toán này có $m = 2$ ràng buộc chính và $n = 3$ biến. Một phương án cực biên có nhiều nhất $m = 2$ thành phần dương, tức là có ít nhất $n - m = 1$ thành phần bằng $0$.
 Vì thế, lần lượt cho $x_1, x_2, x_3$ bằng $0$ ta được:


 Với $x_1 = 0$, hệ phương trình trên cho ta $x_2 = \dfrac{9}{2}; x_3 = 5$.

 Với $x_2 = 0$, hệ phương trình trên vô nghiệm.

Với $x_3 = 0$, hệ phương trình trên cho ta $x_1 = 5; x_2 =\dfrac{9}{2}.$

 Ta nhận được hai phương án: $\left(0, \dfrac{9}{2}, 5\right)^T$ và $\left(5, \dfrac{9}{2}, 0\right)^T$.

Kiểm tra trực tiếp cho thấy hệ 2 véctơ $A^2 = (-2, 2)^T, $\break
 $A^3 = (3, -1)^T$ và
$A^1 = (3, -1)^T, A^2 = (-2, 2)^T$ là độc lập tuyến tính.

Nên cả hai phương án trên đều là các phương án cực biên không suy biến (số thành phần dương bằng $m = 2$).
\end{Soln}
\begin{Soln}{2}
Giải bằng phương pháp hình học. Miền các phương án được xác định bằng đa giác lồi mở
$xABCDy$ với $A(12,0), B(4,2), C(2,6), D(0,18)$. Biểu diễn vectơ pháp tuyến $\vec{n}=(1,2)$ và hàm mục tiêu $L(x,y)=20x+40y$ như hình vẽ.

\begin{figure}[!ht]
\centering
%\input hinh24.tex
\includegraphics[scale =0.7]{hinh24}
\caption{}\label{fig:2}
\end{figure}

Di chuyển đường hàm mục tiêu ngược với $\vec{n}$ tới vị trí giới hạn là điểm $B(4,2)$. Vậy bài toán không có phương án tối ưu.

\end{Soln}
\begin{Soln}{3}
a) Đưa bài toán về dạng chính tắc
$$ f=5x_1+4x_2+6x_3\rightarrow \min$$
\begin{equation*}
\left\{
\begin{array}{r r r r r  r l }
2x_1&+4x_2&+3x_3&-x_4&&&= 48 \\
3x_1&+2x_2&+3x_3&&-x_5&&= 30 \\
x_1&+3x_2&+4x_3&&&-x_6&= 23 \\
\multicolumn{6}{l}{x_j\ge 0,\quad j = 1, 2, 3, 4, 5, 6.}\\
\end{array}
\right.
\end{equation*}
b) Ta biến đổi để xuất hiện biến cơ bản như sau:  Giữ nguyên phương trình thứ nhất, đem nó trừ đi lần lượt các phương trình thứ 2 và 3 cho kết quả mới. Ta có bài toán tương đương
$$ f=5x_1+4x_2+6x_3\rightarrow \min$$
\begin{equation*}
\left\{
\begin{array}{r r r r r  r l }
2x_1&+4x_2&+3x_3&-x_4&&&= 48 \\
-x_1&+2x_2&&-x_4&+x_5&&= 18 \\
x_1&+x_2&-x_3&-x_4&&+x_6&= 25 \\
\multicolumn{6}{l}{x_j\ge 0,\quad j = 1, 2, 3, 4, 5, 6.}\\
\end{array}
\right.
\end{equation*}
Ta đưa về bài toán mở rộng
$$ f=5x_1+4x_2+6x_3\qquad +Mx_7\rightarrow \min$$
\begin{equation*}
\left\{
\begin{array}{r r r r r  r r l }
2x_1&+4x_2&+3x_3&-x_4&&&+x_7&= 48 \\
-x_1&+2x_2&&-x_4&+x_5&&&= 18 \\
x_1&+x_2&-x_3&-x_4&&+x_6&&= 25 \\
\multicolumn{6}{l}{x_j\ge 0,\quad j = 1, 2, 3, 4, 5, 6,7.}\\
\end{array}
\right.
\end{equation*}
Trong đó $x_7$ là biến giả, $M>0$ rất lớn.

\begin{tabular}{| l | l | l | l  l  l  l  l  l  l |}
\hline
$x_B$&$c_B$&PA&$x_1$&$x_2$&$x_3$&$x_4$&$x_5$&$x_6$&$x_7$\\
&&&5&4&6&0&0&0&M\\
\hline
$x_7$&M&48&[2]&4&3&-1&0&0&1\\
$x_5$&0&18&-1&2&0&-1&1&0&0\\
$x_6$&0&25&1&1&-1&-1&0&1&0\\
\hline
f &M&48&2&4&3&-1&0&0&0\\
\cline{2-10}
&&0&-5&-4&-6&0&0&0&0\\
\hline
$x_1$&5&24&1&2&3/2&-1/2&0&0&\\
$x_5$&0&42&0&[4]&3/2&-3/2&1&0&\\
$x_6$&0&1&0&-1&-5/2&-1/2&0&1&\\
\hline
f&&120&0&6&3/2&-5/2&0&0&\\
\hline
$x_1$&5&3&1&0&3/4&1/4&-1/2&0&\\
$x_2$&4&21/2&0&1&3/8&-3/8&1/4&0&\\
$x_6$&0&23/2&0&0&-17/8&-7/8&1/4&1&\\
\hline
f&&57&0&0&-3/4&-1/4&-3/2&0&\\
\hline
\end{tabular}

Bài toán đã cho có phương án tối ưu duy nhất là
$x^*=(3, 21/2, 0)^T$ và $f_{\min}=57$.
\end{Soln}
